\begin{abstract}
In dieser Arbeit wird untersucht, warum sich Kryptowährungen bis 2026 nicht als alltägliches Zahlungsmittel durchgesetzt haben. Dazu werden drei Wirtschaftsräume - El Salvador, die Vereinigten Arabischen Emirate und die Europäische Union – miteinander verglichen. Die Analyse stützt sich auf wissenschaftliche Studien und offizielle Regulierungsdokumente. Die \hyperref[chap:interpretation]{Ergebnisse} zeigen, dass weder eine staatliche Einführung noch eine Regulierung allein ausreicht, sondern dass das Zusammenspiel aus Preisstabilität, gesellschaftlichem Vertrauen und praktischem Nutzen im Alltag entscheidend ist.
\end{abstract}
